\section{Supplementary Problem 2.6}
\subsection{Criterion for time invariance}
Let $y(n)=\mathcal{T}[x](n)$ and let $x_k(n)=x(n-k)$ be a delayed input.
A system is time invariant if
\begin{equation}
  \mathcal{T}[x_k](n)=y(n-k)
\end{equation}
for every input $x$, integer shift $k$, and time $n$. The left side
delays the input before applying the system; the right side delays the
original output. A single counterexample disproves time invariance,
whereas agreement for one example alone does not prove it.

\subsection{Part (a): Squared-time system}
For $\mathcal{T}[x](n)=x(n^2)$, the two paths give
\begin{align}
  \mathcal{T}[x_k](n)&=x_k(n^2)=x(n^2-k),\\
  y(n-k)&=x((n-k)^2).
\end{align}
These are not equal in general. For the pulse input used in part (b),
choose $k=2$ and $n=-2$. Then $x(n^2-k)=x(2)=1$, while
$x((n-k)^2)=x(16)=0$. This counterexample proves that the system
is time varying.

\subsection{Part (b): Explicit pulse example}
The specified input is
\begin{equation}
  x(n)=\begin{cases}1,&0\leq n\leq3,\\0,&\text{otherwise},\end{cases}
  \qquad n\in\mathbb{Z}.
\end{equation}
Use $k=2$ and write $x_2(n)=x(n-2)$,
$y'_2(n)=y(n-2)$, and $y_2(n)=\mathcal{T}[x_2](n)$.
The five signals requested in steps (1)--(5) are determined as follows;
all unlisted samples are zero.
\begin{table}[H]
  \centering
  \caption{Pulse example for the squared-time system}
  \label{tab:p26b}
  \begin{tabular}{lll}
    \toprule
    Signal & Nonzero time indices & Sample values\\
    \midrule
    $x(n)$ & $0,1,2,3$ & $1,1,1,1$\\
    $y(n)=x(n^2)$ & $-1,0,1$ & $1,1,1$\\
    $y'_2(n)=x((n-2)^2)$ & $1,2,3$ & $1,1,1$\\
    $x_2(n)=x(n-2)$ & $2,3,4,5$ & $1,1,1,1$\\
    $y_2(n)=x(n^2-2)$ & $-2,2$ & $1,1$\\
    \bottomrule
  \end{tabular}
\end{table}
For the original output, $0\leq n^2\leq3$ gives $n=-1,0,1$.
For the shifted-input output, $0\leq n^2-2\leq3$ gives
$2\leq n^2\leq5$, whose only integer solutions are $n=\pm2$.
For step (6), the discrepancy is
\begin{equation}
  y_2(n)-y'_2(n)=\delta(n+2)-\delta(n-1)-\delta(n-3).
\end{equation}
The two paths differ, confirming the conclusion in part (a).

\subsection{Part (c): First-difference system}
Now consider $\mathcal{T}[x](n)=x(n)-x(n-1)$. For an arbitrary input
and arbitrary integer delay,
\begin{align}
  \mathcal{T}[x_k](n)&=x(n-k)-x(n-1-k),\\
  y(n-k)&=x(n-k)-x(n-k-1).
\end{align}
The two expressions are identical for all $x,n,k$. Hence this system
is time invariant.

For the same pulse input, the original output has one positive sample
at the leading edge and one negative sample immediately after the pulse:
\begin{equation}
  y(n)=\delta(n)-\delta(n-4).
\end{equation}
Repeating steps (1)--(5) gives the following nonzero samples:
\begin{table}[H]
  \centering
  \caption{Pulse example for the first-difference system}
  \label{tab:p26c}
  \begin{tabular}{lll}
    \toprule
    Signal & Nonzero time indices & Sample values\\
    \midrule
    $x(n)$ & $0,1,2,3$ & $1,1,1,1$\\
    $y(n)$ & $0,4$ & $1,-1$\\
    $y'_2(n)=y(n-2)$ & $2,6$ & $1,-1$\\
    $x_2(n)$ & $2,3,4,5$ & $1,1,1,1$\\
    $y_2(n)=\mathcal{T}[x_2](n)$ & $2,6$ & $1,-1$\\
    \bottomrule
  \end{tabular}
\end{table}
Thus $y_2(n)=y'_2(n)=\delta(n-2)-\delta(n-6)$ and their difference
is zero. This particular agreement illustrates time invariance; the
general algebraic identity above is what proves it.

\subsection{Part (d): Multiplication by time}
For $\mathcal{T}[x](n)=n x(n)$, the two paths are
\begin{align}
  \mathcal{T}[x_k](n)&=n x(n-k),\\
  y(n-k)&=(n-k)x(n-k).
\end{align}
Their difference is $k x(n-k)$, which is nonzero for suitable inputs
and nonzero $k$. Consequently this system is time varying.

For the given pulse and $k=2$, the nonzero samples in steps (1)--(5) are
\begin{table}[H]
  \centering
  \caption{Pulse example for the time-multiplication system}
  \label{tab:p26d}
  \begin{tabular}{lll}
    \toprule
    Signal & Nonzero time indices & Sample values\\
    \midrule
    $x(n)$ & $0,1,2,3$ & $1,1,1,1$\\
    $y(n)=nx(n)$ & $1,2,3$ & $1,2,3$\\
    $y'_2(n)=(n-2)x(n-2)$ & $3,4,5$ & $1,2,3$\\
    $x_2(n)$ & $2,3,4,5$ & $1,1,1,1$\\
    $y_2(n)=nx(n-2)$ & $2,3,4,5$ & $2,3,4,5$\\
    \bottomrule
  \end{tabular}
\end{table}
The difference $y_2(n)-y'_2(n)=2x(n-2)$ equals two at
$n=2,3,4,5$ and zero elsewhere. For example, at $n=2$ the shifted
output is zero while the response to the shifted input is two.

\subsection{Scilab implementation and expected results}
The script uses a direct evaluator for the pulse to avoid truncating
shifted signals at vector boundaries. It creates one figure per system,
each containing the five requested signals and a sixth plot of
$\mathcal{T}[x_2](n)-y(n-2)$. The interval $-4\leq n\leq8$ contains
all nonzero samples for these examples. The maximum absolute
discrepancies are expected to be $1$, $0$, and $2$ for parts (b), (c),
and (d), respectively.
\lstinputlisting[caption={Scilab implementation of supplementary Problem 2.6},label={lst:p26}]{../2_6.sci}

\subsection{Execution evidence}
\resultfigure{images/problem26-b.png}
  {Pulse experiment for the squared-time system in Problem 2.6(b)}
  {fig:p26b}
\resultfigure{images/problem26-c.png}
  {Pulse experiment for the first-difference system in Problem 2.6(c)}
  {fig:p26c}
\resultfigure{images/problem26-d.png}
  {Pulse experiment for the time-multiplication system in Problem 2.6(d)}
  {fig:p26d}
